\section*{Limitations}
\label{sec:limitation}

\paragraph{Corpus Imbalance and Domain Coverage.}
Although our corpus encompasses 18,438 pairs across 26 tasks, the distribution remains inherently skewed.
Mainstream paradigms like Classification and Regression dominate the dataset, whereas niche scientific tasks (e.g., Audio Classification, Tabular Grading) are represented by significantly smaller sample sizes.
Consequently, while the model demonstrates strong general capabilities, its reliability in extremely low-resource or highly specialized scientific domains may vary, and the current evaluation may not fully reflect the challenges of these long-tail scenarios.
Furthermore, the Verified Data Report currently relies on metadata for unstructured domains like CV and NLP, leaving the integration of multimodal data analysis agents for deeper semantic profiling to future work.


\paragraph{Agent Framework Implementation.}
To validate the model's utility, we prioritized stability, instantiating \ours with a conservative \textit{Predict-then-Verify} loop.
This design alternates strictly between singular prediction and execution, barely scratching the surface of potential inference-time strategies.
Specifically, we have not exhaustively explored advanced architectural variants or hyperparameter configurations within this paradigm, implying that the current implementation has not yet been pushed to its optimal limit.
Therefore, the reported performance likely represents a lower bound of the framework's capability.
Beyond this specific instantiation, we identify the framework's broader potential as a scalable \textit{Reward Model}.
By providing dense, execution-free feedback, it paves the way for accelerating Reinforcement Learning rollouts and serves as a plug-and-play optimization module adaptable to diverse agent frameworks.
